\chapter{Quasi-likelihood and moment estimators}
\label{ch:quasi-likelihood-moments}

The paper's headline estimator maximises the exact marginal likelihood of the noised
data by expectation--maximisation around the tree-exact recursion
(\Cref{ch:maximum-likelihood-em,ch:em-through-noise}). Two simpler estimators appear
beside it. The first is the \emph{Gaussian quasi-likelihood}: the Gaussian
log-likelihood written down for data that are not Gaussian. The paper needs it in one
sentence of Section~2 (a Gaussian-innovation kernel fitted by maximum likelihood is
consistent for $(\corr,\ivar)$ and lands on the LMMSE floor) and uses that sentence as
an apparatus check in every capacity table and in the fifteen cells of the regime
sweep. The second is the \emph{method of moments}: two closed-form statistics, one for
the correlation and one for the fourth cumulant, which the paper places against EM as
the rate-correct, inefficient competitor at high noise (Section~4 of the paper,
experiment E19). This chapter builds both from scratch, derives the two properties the
paper asserts (consistency of the quasi-MLE; unbiasedness of the moment estimators),
computes their variances and their efficiency relative to maximum likelihood, and
records how E19 measured them.

\section{The Gaussian quasi-likelihood of noised data}
\label{sec:ql-def}

\subsection{Likelihood of the wrong model}

Fix one noise level $t$ and write $\alpha=e^{-t}$, $\Dt=1-\alpha^2$. A noised sequence
is $\vx=\alpha\va+\sqrt{\Dt}\,\vz$ with $\va\sim\Pz$ and $\vz\sim\normal{0}{\Id}$.
Whatever the innovation law, the chain of the paper has
$\E\va=0$ and $\Cov(\va)=\Sig_0$ with $(\Sig_0)_{ij}=\corr^{|i-j|}$, so
\begin{equation}
  \E\vx=0,\qquad \Cov(\vx)=\alpha^2\Sig_0+\Dt\Id=:\Sig_t .
  \label{eq:ql-cov}
\end{equation}
Only the first two moments of $\Pz$ enter \cref{eq:ql-cov}; this is the whole reason the
Gaussian quasi-likelihood works.

\begin{definition}[Gaussian quasi-likelihood]
\label{def:ql}
Let $\Sig_t(\params)=e^{-2t}\Sig_0(\params)+\Dt\Id$ be the covariance that the Gaussian
chain with parameters $\params=(\corr,\ivar)$ would give to a noised sequence at level $t$.
The \term{Gaussian quasi-likelihood} of $\nseq$ noised sequences $\vx^1,\dots,\vx^{\nseq}$
observed at level $t$ is the Gaussian log-likelihood
\begin{equation}
  \ell_{\mathrm G}(\params)
  = -\frac{1}{2}\sum_{\mu=1}^{\nseq}\Bigl[\log\det\Sig_t(\params)
    + \vx^{\mu\top}\Sig_t(\params)^{-1}\vx^{\mu}\Bigr] - \frac{\nseq\nsites}{2}\log 2\pi ,
  \label{eq:ql}
\end{equation}
evaluated on data whose true law is not Gaussian. Its maximiser is the
\term{quasi-maximum-likelihood estimator} (quasi-MLE) $\hat\params_{\mathrm G}$.
\end{definition}

The word \emph{quasi} records that \cref{eq:ql} is not the log-density of the data. The
true log-density of $\vx$ is $\log\Pt(\vx)$, a non-Gaussian function whose evaluation
needs the forward--backward sweep of \Cref{ch:factor-graphs-bp}. The Gaussian
quasi-likelihood is what one obtains by pretending that the innovation law is Gaussian,
and \cref{eq:ql} is then exactly the marginal likelihood of the Gaussian chain, since the
Gaussian chain noised by a Gaussian channel is Gaussian with covariance $\Sig_t(\params)$.

\subsection{It depends on the data only through the sample covariance}

Let $\mat{S}=\nseq^{-1}\sum_\mu\vx^\mu\vx^{\mu\top}$ be the \term{sample covariance} (about
zero, because the mean is known to be zero). Using
$\vx^\top\mat{A}\vx=\tr(\mat{A}\vx\vx^\top)$ and the linearity of the trace,
\begin{equation}
  \sum_{\mu}\vx^{\mu\top}\Sig_t^{-1}\vx^\mu
  =\tr\Bigl(\Sig_t^{-1}\sum_\mu\vx^\mu\vx^{\mu\top}\Bigr)=\nseq\,\tr\bigl(\Sig_t^{-1}\mat{S}\bigr),
\end{equation}
so
\begin{equation}
  \ell_{\mathrm G}(\params)=-\frac{\nseq}{2}\Bigl[\log\det\Sig_t(\params)+\tr\bigl(\Sig_t(\params)^{-1}\mat{S}\bigr)\Bigr]+\text{const}.
  \label{eq:ql-S}
\end{equation}
This is the statement in the paper's Section~2: \emph{the Gaussian quasi-likelihood of
noised data depends on the data only through their covariance}. Two remarks make it
precise. First, the dependence is through the \emph{sample} covariance $\mat{S}$, a
random matrix; the population statement is about its expectation
$\E\mat{S}=\Sig_t=e^{-2t}\Sig_0+\Dt\Id$, which is \cref{eq:ql-cov}. Second, nothing about
the shape of $\Pz$ beyond its covariance can influence $\hat\params_{\mathrm G}$: two
data sets with the same $\mat{S}$ give the same quasi-MLE, whether their innovations were
Laplace, uniform or Gaussian.

\subsection{The one-view protocol: a sum over levels}

Under the paper's one-view protocol each sequence is observed at one level $t_\mu$ and the
$\nseq$ sequences are split evenly across the twelve levels of the schedule. The
quasi-likelihood is then a sum of twelve terms of the form \cref{eq:ql-S}, one per level,
each with its own sample covariance $\mat{S}_t$ over the $\nseq/12$ sequences at that
level:
\begin{equation}
  \ell_{\mathrm G}(\params)=-\frac{1}{2}\sum_{t}\nseq_t\Bigl[\log\det\Sig_t(\params)+\tr\bigl(\Sig_t(\params)^{-1}\mat{S}_t\bigr)\Bigr]+\text{const},
  \qquad \nseq_t=\nseq/12 .
  \label{eq:ql-levels}
\end{equation}
Every level identifies the same $\params$, so the sum only adds information; the
argument below is given for one level and applies term by term.

\section{Consistency of the quasi-MLE for \texorpdfstring{$(\corr,\ivar)$}{(rho, q)}}
\label{sec:ql-consistency}

\subsection{The population objective is a Kullback--Leibler divergence}

Divide \cref{eq:ql-S} by $\nseq$ and let $\nseq\to\infty$. By the law of large numbers
$\mat{S}\to\Sig_t^{\star}:=e^{-2t}\Sig_0+\Dt\Id$, the true covariance of the noised data,
so the normalised quasi-likelihood converges, uniformly on compact parameter sets, to the
\term{population quasi-likelihood}
\begin{equation}
  \bar\ell(\params)=-\frac12\Bigl[\log\det\Sig_t(\params)+\tr\bigl(\Sig_t(\params)^{-1}\Sig_t^\star\bigr)\Bigr].
  \label{eq:ql-pop}
\end{equation}
Compare this with the Kullback--Leibler divergence between two centred Gaussians
(\Cref{ch:conventions} for the definition of $\KL$):
\begin{equation}
  \KL\bigl(\normal{0}{\Sig^\star}\,\big\|\,\normal{0}{\Sig}\bigr)
  =\frac12\Bigl[\tr(\Sig^{-1}\Sig^\star)-\nsites+\log\det\Sig-\log\det\Sig^\star\Bigr].
  \label{eq:kl-gauss}
\end{equation}
The derivation of \cref{eq:kl-gauss} is short. With $p=\normal{0}{\Sig^\star}$ and
$q=\normal{0}{\Sig}$, $\log(p/q)=\tfrac12[\log\det\Sig-\log\det\Sig^\star]-\tfrac12\vx^\top(\Sig^{\star-1}-\Sig^{-1})\vx$, and
$\E_p[\vx^\top\mat{A}\vx]=\tr(\mat{A}\Sig^\star)$ gives
$\E_p[\vx^\top\Sig^{\star-1}\vx]=\nsites$ and $\E_p[\vx^\top\Sig^{-1}\vx]=\tr(\Sig^{-1}\Sig^\star)$.
Hence
\begin{equation}
  \bar\ell(\params)=-\KL\bigl(\normal{0}{\Sig_t^\star}\,\big\|\,\normal{0}{\Sig_t(\params)}\bigr)+\text{const},
\end{equation}
where the constant does not depend on $\params$. Maximising the population
quasi-likelihood is minimising a Kullback--Leibler divergence between Gaussians, and that
divergence is zero if and only if $\Sig_t(\params)=\Sig_t^\star$, because
$\KL(p\,\|\,q)\ge0$ with equality only at $p=q$ (\Cref{ch:conventions}).

\subsection{Identifiability of \texorpdfstring{$(\corr,\ivar)$}{(rho, q)} from the noised covariance}

It remains to show that $\Sig_t(\params)=\Sig_t^\star$ forces $\params=\params^\star$. Since
$\Sig_t(\params)-\Dt\Id=e^{-2t}\Sig_0(\params)$ and $e^{-2t}>0$, equality of the two noised
covariances is equality of the two clean covariances, $\Sig_0(\params)=\Sig_0(\params^\star)$.
For the chain with $a_1\sim\normal01$ and general $(\corr,\ivar)$ (the two-parameter
family the EM fits), the code's \code{theory.chain\_covariance} gives
\begin{equation}
  (\Sig_0)_{ij}=\corr^{|i-j|}\,v_{\min(i,j)},\qquad
  v_i=\corr^{2(i-1)}+\ivar\,\frac{1-\corr^{2(i-1)}}{1-\corr^2},
  \label{eq:sig0-general}
\end{equation}
so $(\Sig_0)_{11}=1$, $(\Sig_0)_{12}=\corr$ and $(\Sig_0)_{22}=\corr^2+\ivar$: the
$(1,2)$ entry reads off $\corr$ and the $(2,2)$ entry then reads off $\ivar$. Two
different parameter pairs give two different covariances, so $\params$ is identifiable
from $\Sig_0$, hence from $\Sig_t^\star$.

\begin{proposition}[Consistency of the Gaussian quasi-MLE]
\label{prop:qmle}
Let the noised data at level $t$ have a law with mean zero and covariance
$\Sig_t^\star=e^{-2t}\Sig_0(\params^\star)+\Dt\Id$, and let $\hat\params_{\mathrm G}$
maximise \cref{eq:ql} over a compact set containing $\params^\star$ in its interior. Then
$\hat\params_{\mathrm G}\to\params^\star$ in probability as $\nseq\to\infty$, whatever the
innovation law.
\end{proposition}
\begin{proof}
The normalised objective converges uniformly to \cref{eq:ql-pop}, which is continuous
and has the unique maximiser $\params^\star$ by the two facts above; a standard
argument (the argmax of a uniformly convergent sequence of continuous functions converges
to the unique argmax of the limit) finishes the proof.
\end{proof}

This is the special case of White's theory of misspecified maximum likelihood
(\Cref{ch:misspecification-edgeworth}) in which the pseudo-true parameter coincides with
the true one: the model is wrong about the shape of the law but right about its first two
moments, and the Gaussian likelihood only ever asks about those.

\subsection{The quasi-score equations are moment equations}

Differentiate \cref{eq:ql-S} with respect to a parameter $\theta$ using
$\partial\log\det\Sig=\tr(\Sig^{-1}\partial\Sig)$ and
$\partial\Sig^{-1}=-\Sig^{-1}(\partial\Sig)\Sig^{-1}$:
\begin{equation}
  \frac{\partial\ell_{\mathrm G}}{\partial\theta}
  =\frac{\nseq}{2}\,\tr\Bigl(\Sig_t^{-1}\,\partial_\theta\Sig_t\,\Sig_t^{-1}\,(\mat{S}-\Sig_t)\Bigr).
  \label{eq:quasi-score}
\end{equation}
The \term{quasi-score} vanishes when a weighted trace of $\mat{S}-\Sig_t(\params)$
vanishes: the quasi-MLE solves \emph{moment equations}, one per parameter, that ask the
sample covariance to agree with the model covariance in the directions
$\mat{A}_\theta=\Sig_t^{-1}\partial_\theta\Sig_t\Sig_t^{-1}$. That is the sense in which
the quasi-MLE belongs in this chapter beside the method of moments: it is a moment
estimator with a particular, likelihood-derived choice of weights.

\subsection{Variance of the quasi-MLE: the sandwich}
\label{sec:ql-sandwich}

The quasi-MLE is consistent for every innovation law, but its \emph{variance} is not the
Gaussian one when the data are not Gaussian. The standard result for an estimator that
solves $\sum_\mu\psi(\vx^\mu;\params)=0$ (an \term{M-estimator}) is
\begin{equation}
  \sqrt{\nseq}\,(\hat\params-\params^\star)\ \to\ \normal{0}{\mat{J}^{-1}\mat{V}\mat{J}^{-1}},
  \qquad \mat{J}=-\E\,\partial_\params\psi,\quad \mat{V}=\Var\,\psi ,
  \label{eq:sandwich}
\end{equation}
the \term{sandwich covariance}. For the quasi-score $\psi_\theta=\tfrac12\tr(\mat{A}_\theta(\vx\vx^\top-\Sig_t))$
the outer matrices $\mat{J}$ are the Gaussian Fisher information
$\tfrac12\tr(\Sig_t^{-1}\partial_i\Sig_t\Sig_t^{-1}\partial_j\Sig_t)$ of the paper's
Theorem~3, and $\mat{V}$ needs the variance of a quadratic form under the true law.

\begin{lemma}[Variance of a quadratic form]
\label{lem:quadform}
Let $\vx$ have mean zero, covariance $\Sig$ and fourth-cumulant tensor $\kfour[\vx]$
(\Cref{ch:moments-cumulants}), and let $\mat{A}$ be symmetric. Then
\begin{equation}
  \Var\bigl(\vx^\top\mat{A}\vx\bigr)=2\tr(\mat{A}\Sig\mat{A}\Sig)+\sum_{ijkl}A_{ij}A_{kl}\,(\kfour[\vx])_{ijkl}.
\end{equation}
\end{lemma}
\begin{proof}
$\Var(\vx^\top\mat{A}\vx)=\sum_{ijkl}A_{ij}A_{kl}\bigl(\E[x_ix_jx_kx_l]-\Sig_{ij}\Sig_{kl}\bigr)$.
The fourth moment of a zero-mean law is its fourth cumulant plus the three pairings of
second moments (\Cref{ch:moments-cumulants}):
$\E[x_ix_jx_kx_l]=\kfour[ijkl]+\Sig_{ij}\Sig_{kl}+\Sig_{ik}\Sig_{jl}+\Sig_{il}\Sig_{jk}$.
The first pairing cancels against $\Sig_{ij}\Sig_{kl}$; the other two each contract to
$\tr(\mat{A}\Sig\mat{A}\Sig)$ by symmetry of $\mat{A}$.
\end{proof}

For the noised chain, $\kfour[\vx]=e^{-4t}\kfour[\va]$ exactly (the cumulant consequence
of the paper's Theorem~2, \Cref{ch:attenuation-theorems}), so the non-Gaussian part of
$\mat{V}$ is $e^{-4t}$ times a fixed contraction of $\kfour[\va]$, while the Gaussian part
$2\tr(\mat{A}\Sig_t\mat{A}\Sig_t)$ is of order $e^{-4t}$ as well through
$\mat{A}_\theta\propto\partial_\theta\Sig_t=e^{-2t}\partial_\theta\Sig_0$. The ratio of
the two is therefore of order one at every level, not decaying: the quasi-MLE of a
Laplace chain is consistent but has a different (typically larger) variance than the
same estimator on a Gaussian chain, by an amount proportional to the excess kurtosis.
This matters only for error bars; the paper never quotes a variance of the
$\ncomp=1$ arm, so the consistency statement is the one it uses.

\subsection{Why this is an apparatus check}

Put together with Proposition~1 of the paper (\Cref{ch:gaussian-closure-lmmse}), the
chain of implications is:
\begin{enumerate}
\item the $\ncomp=1$ arm of EM--BP fits a Gaussian-innovation kernel by maximum
  likelihood of the noised data; with a Gaussian kernel the exact marginal likelihood that
  EM maximises \emph{is} the Gaussian quasi-likelihood \cref{eq:ql-levels} (up to the grid
  discretisation), so EM's fixed point is the quasi-MLE;
\item by \Cref{prop:qmle} the quasi-MLE converges to the true $(\corr,\ivar)$;
\item by Proposition~1, sum-product with any Gaussian-innovation kernel returns the LMMSE
  estimator of that kernel's covariance; at the true $(\corr,\ivar)$ this is the LMMSE of
  the true covariance, the second-order floor.
\end{enumerate}
So the $\ncomp=1$ row of every capacity table must sit on the LMMSE floor, up to the
$O(\nseq^{-1/2})$ sampling error of step~2 and the checkpoint selection. The paper
reports that it does, at every sample size, at every member of the generalised-Gaussian
family and in all fifteen $(\corr,\nsites)$ cells of the regime sweep (see the paper,
Section~5 and Table~1, and the regime paragraph; E02a, E02b, E13). A failure of that
identity anywhere would indict the grid, the message passing, the M-step or the scoring,
which is why the audit calls it the most reassuring fact about the pipeline.

\section{The method of moments in general}
\label{sec:mom-general}

\begin{definition}[Method of moments]
\label{def:mom}
Let $\vx^1,\dots,\vx^{\nseq}$ be i.i.d.\ with law $P_{\params}$, and let
$g(\vx)\in\Reals^d$ be a vector of statistics whose expectation
$m(\params)=\E_{\params}\,g(\vx)$ is known in closed form and injective in $\params$. The
\term{method-of-moments estimator} is the solution of
\begin{equation}
  \frac{1}{\nseq}\sum_{\mu=1}^{\nseq}g(\vx^\mu)=m(\hat\params),
  \qquad\text{i.e.}\qquad \hat\params=m^{-1}(\bar g).
  \label{eq:mom}
\end{equation}
\end{definition}

Three facts follow from the definition. First, $\bar g$ is an unbiased estimator of
$m(\params)$ for every $\nseq$, so when $m$ is linear the estimator $\hat\params$ is
exactly unbiased; when $m$ is nonlinear the plug-in $m^{-1}(\bar g)$ is unbiased only to
order $\nseq^{-1}$ (the delta method, \Cref{ch:experimental-statistics}). Second, by
the central limit theorem and the delta method,
\begin{equation}
  \Var(\hat\params)\approx\frac{1}{\nseq}\,(\partial m)^{-1}\,\Var(g)\,(\partial m)^{-\top},
  \label{eq:mom-var}
\end{equation}
which is \cref{eq:sandwich} with $\psi=g-m(\params)$ and $\mat{J}=\partial m$. Third,
the estimator is asymptotically efficient, that is, attains the Cram\'er--Rao bound
(\Cref{ch:fisher-cramer-rao}), exactly when the statistic $g$ is an affine function of
the score $\partial_\params\log P_\params$: otherwise the ratio of the Cram\'er--Rao
variance to \cref{eq:mom-var}, the \term{asymptotic relative efficiency}, is strictly
below one.

\paragraph{Moment estimators through the channel.}
The paper's Theorem~2 (Hermite attenuation, \Cref{ch:attenuation-theorems}) is what
makes the method of moments available for noised data without any inference: it says
that for every multi-index $\nu$,
\begin{equation}
  \E_{\Pt}[H_\nu(\vx)]=e^{-|\nu|t}\,\E_{\Pz}[H_\nu(\va)],
\end{equation}
so any Hermite moment of the clean law has a known image under the channel, and the
sample Hermite moment of the noised data, divided by $e^{-|\nu|t}$, is an unbiased
estimator of the clean one. The two estimators below are the simplest instances, at
degree two and at degree four.

\section{The paper's two moment estimators}
\label{sec:two-estimators}

Throughout this section the data are $\nseq$ sequences observed at one common level
$t$, as in E06 and E19 (not the twelve-level split), and $\alpha=e^{-t}$.

\subsection{The lag-one product for the correlation}

\begin{definition}[Lag-one moment estimator]
\label{def:rho-mom}
\begin{equation}
  \hat\corr_{\mathrm{mom}}=\frac{1}{\alpha^2}\cdot\frac{1}{\nseq(\nsites-1)}\sum_{\mu=1}^{\nseq}\sum_{i=1}^{\nsites-1}x^\mu_i\,x^\mu_{i+1}.
  \label{eq:rho-mom}
\end{equation}
\end{definition}

\begin{proposition}[Unbiasedness]
\label{prop:rho-unbiased}
For the chain of the paper, $\E\,\hat\corr_{\mathrm{mom}}=\corr$ exactly, at every $t$,
every $\nseq$ and every innovation law of mean zero and variance $\ivar=1-\corr^2$.
\end{proposition}
\begin{proof}
$x_ix_{i+1}=(\alpha a_i+\sqrt{\Dt}z_i)(\alpha a_{i+1}+\sqrt{\Dt}z_{i+1})$. Expanding, the
cross terms have a factor $\E z_i=0$ or $\E z_{i+1}=0$, and $\E[z_iz_{i+1}]=0$ because the
channel noise is independent across sites; so $\E[x_ix_{i+1}]=\alpha^2\E[a_ia_{i+1}]$.
For the chain, $a_{i+1}=\corr a_i+\varepsilon_{i+1}$ with $\varepsilon_{i+1}$ independent
of $a_i$ and of mean zero, so $\E[a_ia_{i+1}]=\corr\E[a_i^2]=\corr\Var(a_i)$. The
variance choice $\ivar=1-\corr^2$ with $a_1\sim\normal01$ gives $\Var(a_i)=1$ at every
site (induction: $\Var(a_{i+1})=\corr^2\Var(a_i)+\ivar=\corr^2+1-\corr^2=1$). Hence
$\E[x_ix_{i+1}]=\alpha^2\corr$ for every $i$, and dividing by $\alpha^2$ and averaging
gives $\corr$.
\end{proof}

The statistic $x_ix_{i+1}$ is, up to the constant $\alpha^2$, the degree-two Hermite
moment with multi-index $\nu=e_i+e_{i+1}$, that is, the expectation of
$\Hen{1}(x_i)\,\Hen{1}(x_{i+1})$; \Cref{prop:rho-unbiased} is Theorem~2 at that
multi-index. The estimator is a linear function of the data products, so it is
unbiased exactly, not only asymptotically.

\paragraph{Variance and the $e^{2t}$ growth.}
Write $T=\frac{1}{\nseq(\nsites-1)}\sum_{\mu,i}x^\mu_ix^\mu_{i+1}$, so
$\hat\corr_{\mathrm{mom}}=\alpha^{-2}T$ and $\Var(\hat\corr_{\mathrm{mom}})=e^{4t}\Var(T)$.
Sequences are independent, so $\Var(T)=\nseq^{-1}\Var(\bar p)$ with
$\bar p=(\nsites-1)^{-1}\sum_ix_ix_{i+1}$ the within-sequence average of the products.
For a Gaussian chain Isserlis' theorem (\Cref{ch:gaussian-hermite}) gives, with
$c_k=\Cov(x_i,x_{i+k})=\alpha^2\corr^{k}$ for $k\ge1$ and $c_0=1$,
\begin{equation}
  \Var(x_ix_{i+1})=\E[x_i^2x_{i+1}^2]-c_1^2=(1+2c_1^2)-c_1^2=1+c_1^2,
  \qquad
  \Cov(x_ix_{i+1},x_{i+k}x_{i+k+1})=c_k^2+c_{k-1}c_{k+1},
\end{equation}
and for a non-Gaussian chain each of these acquires a fourth-cumulant term of order
$e^{-4t}\kfour$ (\Cref{lem:quadform}). All the covariance terms are $O(\alpha^4)$ or
$O(\alpha^2)$ while the leading $1$ is not, so as $t$ grows
\begin{equation}
  \Var(\hat\corr_{\mathrm{mom}})=\frac{e^{4t}}{\nseq(\nsites-1)}\bigl(1+O(e^{-2t})\bigr),
  \qquad
  \mathrm{RMSE}(\hat\corr_{\mathrm{mom}})=\frac{e^{2t}}{\sqrt{\nseq(\nsites-1)}}\bigl(1+O(e^{-2t})\bigr).
  \label{eq:rho-mom-var}
\end{equation}
The error of the degree-two moment estimator grows as $e^{2t}$: the amplitude of the
degree-two signal shrinks as $e^{-2t}$ while the noise in the product stays of order one.
This is the first of the two slopes E19 fits (see the paper, Section~4 and Appendix A5).

\subsection{The per-site fourth cumulant for the shape}

The shape coordinate the paper tracks is the excess kurtosis of the innovation,
$\kfour(\varepsilon)/\ivar^2$, which is $+3$ for the Laplace law. To estimate it from
noised data one needs the fourth cumulant of a site of the chain and its image under the
channel.

\begin{lemma}[Fourth cumulant of a chain site]
\label{lem:kappa4-site}
For the chain $a_1\sim\normal01$, $a_i=\corr a_{i-1}+\varepsilon_i$ ($i\ge2$) with
i.i.d.\ innovations of fourth cumulant $\kfour(\varepsilon)$,
\begin{equation}
  \kfour(a_i)=\kfour(\varepsilon)\sum_{m=2}^{i}\corr^{4(i-m)}
  =\kfour(\varepsilon)\,\frac{1-\corr^{4(i-1)}}{1-\corr^4}=:\kfour(\varepsilon)\,c_i(\corr).
  \label{eq:kappa4-site}
\end{equation}
\end{lemma}
\begin{proof}
Unrolling the recursion, $a_i=\corr^{i-1}a_1+\sum_{m=2}^{i}\corr^{i-m}\varepsilon_m$, a
sum of independent terms. Cumulants are additive over independent summands and
homogeneous of degree four under scaling (\Cref{ch:moments-cumulants}):
$\kfour(\lambda Y)=\lambda^4\kfour(Y)$. The Gaussian $a_1$ has $\kfour(a_1)=0$, so
$\kfour(a_i)=\sum_{m=2}^{i}\corr^{4(i-m)}\kfour(\varepsilon)$, a geometric sum with ratio
$\corr^4$ and $i-1$ terms.
\end{proof}

Under the channel, $x_i=\alpha a_i+\sqrt{\Dt}z_i$ with $z_i$ Gaussian and independent of
$a_i$, so by the same two properties $\kfour(x_i)=\alpha^4\kfour(a_i)$: this is the
cumulant form of Theorem~2 at degree four. Since $x_i$ has mean zero and variance one at
every $t$, $\kfour(x_i)=\E x_i^4-3(\E x_i^2)^2=\E x_i^4-3$.

\begin{definition}[Per-site fourth-cumulant estimator]
\label{def:k4-mom}
With $\hat m_{2,i}=\nseq^{-1}\sum_\mu(x^\mu_i)^2$ and $\hat m_{4,i}=\nseq^{-1}\sum_\mu(x^\mu_i)^4$,
let $\hat k_{4,i}=\hat m_{4,i}-3\hat m_{2,i}^2$ be the sample fourth cumulant at site
$i$. The estimator of the innovation's fourth cumulant, and of its excess kurtosis, is
\begin{equation}
  \hat\kfour(\varepsilon)=\frac{\sum_{i=1}^{\nsites}\hat k_{4,i}}{\alpha^4\sum_{i=1}^{\nsites}c_i(\hat\corr)},
  \qquad
  \widehat{\text{exc.\ kurt.}}=\frac{\hat\kfour(\varepsilon)}{\hat\ivar^{\,2}},\quad
  \hat\ivar=1-\hat\corr^2 ,
  \label{eq:k4-mom}
\end{equation}
with $\hat\corr=\hat\corr_{\mathrm{mom}}$ from \cref{eq:rho-mom} (E19 also records the
version with the true $\corr$ substituted, to separate the two sources of error).
\end{definition}

\begin{proposition}[Unbiasedness at the population level]
\label{prop:k4-unbiased}
Replace the sample cumulants in \cref{eq:k4-mom} by their population values and
$\hat\corr$ by $\corr$. Then the right-hand side equals $\kfour(\varepsilon)$ exactly, at
every $t$.
\end{proposition}
\begin{proof}
$\sum_i\kfour(x_i)=\alpha^4\sum_i\kfour(a_i)=\alpha^4\kfour(\varepsilon)\sum_ic_i(\corr)$
by \Cref{lem:kappa4-site} and the cumulant scaling; divide.
\end{proof}

This is the sense in which the paper's Section~4 says both estimators are ``unbiased by
Theorem~2''. It is a statement about the population identity, and it is exact. The
finite-sample estimator, however, is a nonlinear function of sample moments, and its
bias deserves its own paragraph, because the channel amplifies it.

\paragraph{Finite-sample bias of the plug-in cumulant.}
The sample fourth cumulant $\hat k_4=\hat m_4-3\hat m_2^2$ is not unbiased: $\hat m_4$
is unbiased for $m_4$, but $\E[\hat m_2^2]=(\E\hat m_2)^2+\Var(\hat m_2)=m_2^2+(m_4-m_2^2)/\nseq$,
so
\begin{equation}
  \E\hat k_4=m_4-3m_2^2-\frac{3(m_4-m_2^2)}{\nseq}
  =\kfour\Bigl(1-\frac3\nseq\Bigr)-\frac{6m_2^2}{\nseq}.
  \label{eq:k4-bias}
\end{equation}
At unit variance the bias is $-(3\kfour(x_i)+6)/\nseq$, of which the $-6/\nseq$ term does
not decay with $t$. Summed over $\nsites$ sites and divided by $\alpha^4\sum_ic_i$ and
$\ivar^2$, the bias in excess-kurtosis units is about
\begin{equation}
  -\frac{6\nsites}{\nseq\,\ivar^2\,e^{-4t}\sum_ic_i(\corr)}
  \approx-\frac{6\,(1-\corr^4)}{\nseq\,\ivar^2}\,e^{4t}
  \qquad(\text{large }\nsites),
  \label{eq:k4-bias-amplified}
\end{equation}
which grows as $e^{4t}$: the $O(1/\nseq)$ bias of a moment estimator is multiplied by
the same $e^{4t}$ that multiplies its standard deviation. At $\nseq=512$, $\corr=0.85$
and $\ivar=1-\corr^2=0.2775$ this is about $-0.4$ at $t=0.4$ and about $-4$ at $t=1$,
that is, of the same order as the coordinate's whole range ($3$) at the level where the
paper's own reference bound crosses that range. The remedy is classical: Fisher's
\term{k-statistic} $k_4=\nseq^2[(\nseq+1)\hat m_4-3(\nseq-1)\hat m_2^2]/[(\nseq-1)(\nseq-2)(\nseq-3)]$
is exactly unbiased for $\kfour$ at every $\nseq$. E19 uses the plug-in form, so its
$\hat\kfour$ is unbiased only up to $O(e^{4t}/\nseq)$; the script's docstring says
``unbiased up to $O(1/M)$'', and the paper's Appendix A5 reports the measured bias at
$t=0.4$ (see the paper's macro for that level). A second, smaller $O(1/\nseq)$ bias
comes from dividing by $c_i(\hat\corr)$ and $\hat\ivar^2$, nonlinear functions of the
estimated $\hat\corr$.

\paragraph{Variance and the $e^{4t}$ growth.}
Each $\hat k_{4,i}$ has variance of order $\nseq^{-1}\Var(x_i^4)=O(\nseq^{-1})$ with a
leading constant that does not decay in $t$ (for a standard Gaussian $\Var(x^4)=96$),
while the signal $\kfour(x_i)=\alpha^4\kfour(a_i)$ decays as $e^{-4t}$. Dividing by
$\alpha^4$ therefore multiplies the standard deviation by $e^{4t}$:
\begin{equation}
  \mathrm{RMSE}\bigl(\hat\kfour(\varepsilon)\bigr)\propto\frac{e^{4t}}{\sqrt{\nseq}}
  \qquad(t\to\infty).
\end{equation}
The two slopes, $2$ for the degree-two coordinate and $4$ for the degree-four one, are
the ``three clocks'' of the paper's Section~3 read in the sample-size clock, and they
are what the paper's Appendix A5 reports E19 measuring on $t\ge1$ (the log-slopes are
quoted there as macros).

\paragraph{No collapse.}
EM at high noise collapses the shape estimate to a Gaussian innovation (fitted excess
kurtosis near zero), because the likelihood surface is flat in the shape direction and
the validation-selected early stopping shrinks toward the initialisation
(\Cref{ch:em-through-noise}; the paper's Section~4). The moment estimator cannot do
this: it is a polynomial in the data with no boundary, so its error grows without
saturating and its sign is unconstrained. The paper uses this contrast in one sentence
(``never collapse''), and it is the reason the asymptotic slopes are visible in E19 where
they are not in E06.

\section{Relative efficiency against maximum likelihood}
\label{sec:mom-efficiency}

\subsection{The reference}

\Cref{ch:fisher-cramer-rao} defines the Fisher information and the Cram\'er--Rao bound;
\Cref{ch:attenuation-theorems} gives the paper's Theorem~3, which for one noised
sequence at level $t$ reads
\begin{equation}
  I_t(\corr)=\tfrac12\Frob{\partial_\corr\Sig_0}^2e^{-4t}\bigl(1+O(e^{-2t})\bigr),
  \qquad
  I_t(\beta)=\tfrac1{24}\Frob{\partial_\beta\kfour}^2e^{-8t}\bigl(1+O(e^{-2t})\bigr),
\end{equation}
for a covariance coordinate and a shape coordinate respectively. The reference RMSE for
an unbiased estimator from $\nseq$ sequences is $1/\sqrt{\nseq I_t}$, which grows as
$e^{2t}$ and $e^{4t}$: the moment estimators have the right \emph{rates} by
construction. What separates them from maximum likelihood is the \emph{constant}.

\subsection{The lag-one estimator against the covariance information}

For the Gaussian chain the information about $\corr$ is exact at every $t$
(the paper's Theorem~3, last sentence), and its $t\to\infty$ form is
$e^{-4t}\cdot\tfrac12\Frob{\partial_\corr\Sig_0}^2$. With
$(\Sig_0)_{ij}=\corr^{|i-j|}$, the derivative has entries
$|i-j|\corr^{|i-j|-1}$ off the diagonal and zero on it, so
\begin{equation}
  \tfrac12\Frob{\partial_\corr\Sig_0}^2=\tfrac12\sum_{i\ne j}(|i-j|)^2\corr^{2(|i-j|-1)}
  =\sum_{k=1}^{\nsites-1}(\nsites-k)\,k^2\corr^{2(k-1)}.
  \label{eq:info-rho-constant}
\end{equation}
Comparing with \cref{eq:rho-mom-var}, the asymptotic relative efficiency of the lag-one
estimator is
\begin{equation}
  \mathrm{ARE}(\hat\corr_{\mathrm{mom}})=\lim_{t\to\infty}\frac{1/(\nseq I_t(\corr))}{\Var(\hat\corr_{\mathrm{mom}})}
  =\frac{\nsites-1}{\sum_{k=1}^{\nsites-1}(\nsites-k)k^2\corr^{2(k-1)}}\ \le\ 1,
  \label{eq:are-rho}
\end{equation}
with equality if and only if $\corr=0$. For a long chain the per-site sum converges,
$\sum_{k\ge1}k^2r^{k-1}=(1+r)/(1-r)^3$ with $r=\corr^2$, so
\begin{equation}
  \mathrm{ARE}(\hat\corr_{\mathrm{mom}})\ \xrightarrow[\nsites\to\infty]{}\ \frac{(1-\corr^2)^3}{1+\corr^2}.
  \label{eq:are-rho-long}
\end{equation}

The reason is visible in the score. At high noise the likelihood score of $\corr$ is, to
leading order, the degree-two Hermite statistic
$\tfrac12\alpha^2\langle\partial_\corr\Sig_0,{:}\vx\vx^\top{:}\rangle=\tfrac12\alpha^2\sum_{i\ne j}|i-j|\corr^{|i-j|-1}x_ix_j$
(the paper's Appendix A1, ``information and the general order''): it uses \emph{every}
lag $k$, weighted by $k\corr^{k-1}$, because a change of $\corr$ moves the whole
covariance. The lag-one product uses lag one only. When $\corr=0$ the weights vanish for
$k\ge2$ and the two statistics coincide; when $\corr$ is large the longer lags carry most
of the information and the lag-one estimator discards it.

\begin{example}[Efficiency of the lag-one estimator, by hand]
\label{ex:are}
\emph{Independent sites.} At $\corr=0$ the chain is i.i.d.\ $\normal01$, so
$\Sig_0=\Id$, $\Sig_t=\Id$ at every $t$, and the noised sites $x_i$ are i.i.d.\
$\normal01$. The exact Gaussian information is
$I_t(\corr)=\tfrac12\tr(\Sig_t^{-1}\partial\Sig_t\Sig_t^{-1}\partial\Sig_t)=\tfrac12e^{-4t}\Frob{\partial_\corr\Sig_0}^2=e^{-4t}(\nsites-1)$,
since $\partial_\corr\Sig_0$ at $\corr=0$ has ones on the two first off-diagonals and
$2(\nsites-1)$ nonzero entries. On the estimator side, the products $x_ix_{i+1}$ and
$x_{i+1}x_{i+2}$ are uncorrelated ($\E[x_ix_{i+1}^2x_{i+2}]=\E x_i\cdot\E x_{i+1}^2\cdot\E x_{i+2}=0$)
and each has variance $\E x_i^2\,\E x_{i+1}^2=1$, so
$\Var(\hat\corr_{\mathrm{mom}})=e^{4t}/(\nseq(\nsites-1))=1/(\nseq I_t)$ exactly, at every
$t$: the lag-one estimator is efficient. \emph{The frozen correlation.} At $\corr=0.85$,
$\corr^2=0.7225$, $1-\corr^2=0.2775$, $(1-\corr^2)^3=0.02137$ and
$1+\corr^2=1.7225$, so by \cref{eq:are-rho-long} the long-chain efficiency is
$0.02137/1.7225\approx0.0124$: the lag-one product uses about one eightieth of the
covariance information, and its RMSE at high noise is about $\sqrt{80.6}\approx9$ times
the Cram\'er--Rao reference. (At $\nsites=32$ the finite sum in \cref{eq:are-rho} is a
little smaller than the infinite one, so the factor is a little below nine.) The exact
Gaussian information includes the effect of the nuisance $\ivar$ only through the
efficient information, which changes the constant by an $O(1)$ factor
(\Cref{ch:attenuation-theorems}); the order of magnitude is the point.
\end{example}

\subsection{The per-site cumulant against the shape information}

The same argument applies at degree four. The efficient shape score at high noise is
\begin{equation}
  \frac{\alpha^4}{24}\sum_{ijkl}\partial_\beta\kfour[ijkl]\;{:}x_ix_jx_kx_l{:}\,,
\end{equation}
a contraction of the \emph{whole} fourth-cumulant tensor with the degree-four Hermite polynomials,
including every off-diagonal entry $\kfour[ijkl]$ with $i,j,k,l$ not all equal. The
per-site estimator uses only the diagonal entries $\kfour[iiii]$. For the chain the
off-diagonal entries are
\begin{equation}
  \kfour[ijkl]=\kfour(\varepsilon)\sum_{m=2}^{\min(i,j,k,l)}\corr^{\,i+j+k+l-4m}
\end{equation}
(\Cref{ch:moments-cumulants}; \code{theory.chain\_kappa4}), which are comparable to the
diagonal ones when $\corr$ is close to one, so again the moment estimator discards most
of the information at the frozen correlation. Its constant is therefore far from the
Cram\'er--Rao one, as the paper's Section~4 reports (the ratios are quoted there as
macros); its rate is right.

\begin{remark}[On ``the two estimators coincide at high noise'']
\label{rem:coincide}
The paper's Section~4 says that at high noise the moment and likelihood estimators
coincide, ``since the likelihood score of a degree-$m$ coordinate is itself the
degree-$m$ Hermite statistic''. The second clause is the general-order theorem of
\Cref{ch:attenuation-theorems} and is correct. The first clause is correct for the
moment estimator built from that \emph{full} degree-$m$ statistic (all lags at degree
two; the full tensor at degree four). The paper's two estimators use the lag-one
product and the per-site cumulant, which are proper sub-statistics of it; their
efficiency relative to maximum likelihood tends to the constant \cref{eq:are-rho},
not to one, unless $\corr=0$. What does coincide at high noise is the \emph{rate}, and
in a sample of fixed size both estimators are useless there anyway, which is the
practical content of the sentence.
\end{remark}

\section{How E19 measured them}
\label{sec:e19}

E19 (\code{experiments/e19\_moment\_estimators.py}) is a deterministic re-analysis of
E06's data. It regenerates E06's chains and noise draws from their random-stream tags
(\code{e06-train}, \code{e06-noise}; \Cref{ch:protocol} explains the tag mechanism), so
every row is paired by seed and level with the EM--BP row of
\code{outputs/E06/levels.csv}. The design is: the Laplace chain at the frozen
$\corr=0.85$ and $\nsites=32$; $\nseq=512$ sequences per cell, all noised at one level
$t\in\{0.05,0.1,0.2,0.4,0.7,1.0,1.4,2.0\}$; sixteen seeds. Per cell the script computes
$\hat\corr_{\mathrm{mom}}$ of \cref{eq:rho-mom}, the per-site sample cumulants, and
$\hat\kfour(\varepsilon)/\hat\ivar^2$ of \cref{eq:k4-mom} with $\hat\corr$ and,
separately, with the true $\corr$ (column \code{fit\_kurt\_oracle\_rho}). The
predeclared analysis (registry \code{experiments/registry/E19.json}) is the RMSE,
dispersion and bias over the sixteen seeds per level, paired with E06's EM--BP fits, and
the ordinary-least-squares log-slope of the RMSE against $t$ on the levels $t\ge1$,
compared with the predicted $2$ and $4$. The numbers reach the manuscript through
\code{tools/make\_audit\_numbers.py} as the macros \code{audMom\ldots} (RMSE per level,
bias at each level, the two slopes, and the ratio of the moment RMSE to the EM--BP RMSE
over the informative levels $t\le0.7$); the paper quotes them in Section~4 (the learning
paragraph) and in Appendix A5 (the E19 paragraph). Nothing about the moment estimators
is selected: there is no ladder, no validation bundle and no checkpoint, which is what
makes them the natural control for the selection-dependent EM fits.

Two readings the paper draws from E19, in the order it draws them: below the crossover
the moment estimators are several times worse than EM--BP for the correlation and for
the shape, which is the value of exact inference (\Cref{ex:are} gives the analytic
reason for the correlation); above it, where the EM shape estimate has collapsed, the
moment estimators keep growing at the asymptotic rates, which is the corollary's
prediction made visible on an estimator that cannot saturate.

\where{\begingroup\sloppy The paper, Section~2 (the quasi-likelihood consequence of Proposition~1),
Section~4 (the moment estimators against EM--BP) and Appendix A5 (E19); the
$\ncomp=1$ rows of Table~1, of the family figure and of the regime sweep as the
apparatus check (E02a, E02b, E13). Code:
\code{experiments/e19\_moment\_estimators.py},
\code{experiments/e06\_learning\_through\_noise.py},
\code{tools/make\_audit\_numbers.py},
\code{src/markov\_score/theory.py} (\code{chain\_covariance}, \code{chain\_kappa4},
\code{correlation\_information\_constant}),
\code{src/markov\_score/kernels.py} (the Gaussian and mixture M-steps),
\code{experiments/e13\_regime\_sweep.py}.\par\endgroup}

\paragraph{References.}
H.~White, 1982, Maximum likelihood estimation of misspecified models, \emph{Econometrica} 50, 1--25.
P.~J.~Huber, 1967, The behavior of maximum likelihood estimates under nonstandard conditions, \emph{Proceedings of the Fifth Berkeley Symposium}, 221--233.
A.~W.~van der Vaart, 1998, \emph{Asymptotic Statistics}, Cambridge University Press (chapters 4 and 5: method of moments and M-estimators).
R.~A.~Fisher, 1930, Moments and product moments of sampling distributions, \emph{Proceedings of the London Mathematical Society} 30, 199--238 (the k-statistics).
A.~Balanov, T.~Bendory and D.~Edidin, 2026, The generalized method of moments is (almost) statistically efficient in low-SNR Gaussian latent-variable models, arXiv:2605.30095.
D.~Hsu, S.~M.~Kakade and T.~Zhang, 2012, A spectral algorithm for learning hidden Markov models, \emph{Journal of Computer and System Sciences} 78, 1460--1480.
